\documentclass[11pt,a4paper]{article}

\usepackage[T1]{fontenc}
\usepackage[utf8]{inputenc}
\usepackage[french]{babel}
\usepackage{lmodern}
\usepackage{amsmath,amssymb,amsthm}
\usepackage[a4paper,margin=2.55cm]{geometry}
\usepackage{microtype}
\usepackage{booktabs,array}
\usepackage{fancyvrb}
\usepackage[hidelinks]{hyperref}
\usepackage{enumitem}
\usepackage{xurl}
\setlist[itemize]{leftmargin=*,topsep=3pt,itemsep=2pt}
\setlength{\parskip}{4pt}
\setlength{\parindent}{0pt}
\emergencystretch=2em

\newcommand{\Truth}{\operatorname{Truth}}
\newcommand{\Proof}{\operatorname{Proof}}
\newcommand{\ExactRealization}{\operatorname{ExactRealization}}
\newcommand{\FutureContract}{\operatorname{FutureContract}}
\newcommand{\FutureEquivalent}{\operatorname{FutureEquivalent}}
\newcommand{\Outcome}{\operatorname{Outcome}}
\newcommand{\Empty}{\mathsf{Empty}}
% Les identifiants longs doivent pouvoir se couper aux caractères techniques
% afin de respecter la largeur du texte. Une \mbox ici provoquerait des
% débordements dans les marges pour les hashes et les noms de claims.
\newcommand{\code}[1]{\nolinkurl{#1}}
\newcommand{\Source}[2]{\href{https://github.com/JohnDoe-collab-stack/relational-perimeter/blob/1e5aeda745559d62a3819bb3213da7fce3c31ad1/#1}{\nolinkurl{#2}}}
\newcommand{\SourceAt}[3]{\href{https://github.com/JohnDoe-collab-stack/relational-perimeter/blob/1e5aeda745559d62a3819bb3213da7fce3c31ad1/#1\#L#2}{\nolinkurl{#3}}}
\newcommand{\MainAxis}{%
  \begin{gathered}
  \text{exécution certifiée}\\
  \downarrow\\
  \bigl(C(e),\Gamma_e,\sigma_e\bigr)\\
  \downarrow\\
  \text{obligation universelle}\\
  \text{sur toute réalisation exacte de }C(e)
  \end{gathered}}
\newcommand{\FormalRevision}{1e5aeda745559d62a3819bb3213da7fce3c31ad1}

\newtheorem{theoreme}{Théorème}[section]
\newtheorem{proposition}[theoreme]{Proposition}
\newtheorem{lemme}[theoreme]{Lemme}
\theoremstyle{remark}
\newtheorem{remarque}[theoreme]{Remarque}

\title{De l'exécution gödelienne aux obligations universelles\\
  de réalisation exacte\\[5pt]
  \large Constitution opérationnelle des contrats, séparation et progression finie}
\author{Formalisation \code{relational-perimeter}}
\date{Révision formelle \code{1e5aeda74555}\\
  \small Version documentaire révisée du 23 octobre 2026}

\begin{document}
\maketitle
\vspace{-1.4em}

\begin{abstract}
Dans la théorie formelle plus large des contrats de futurs, du regroupement
opérationnel, de la séparation finie, de la préservation dirigée des readouts
et de la réalisation exacte, le résultat présenté ici est la construction
spécifiquement gödelienne qui dérive une obligation universelle d'une
exécution certifiée et renouvelle cette construction propre à chaque
exécution par une progression finie des régimes de preuve.

Nous présentons un résultat dont l'objet n'est pas seulement la non-prouvabilité
produite par la diagonalisation gödelienne, mais l'obligation que son
\emph{exécution certifiée} impose à toute réalisation exacte d'un contrat
constitué à partir de cette même exécution. À chaque étage arithmétique
certifié, une production exécutée fournit une phrase $G$ et des témoins de
$\Truth(G)$ et de $\Proof_T(G)\to\Empty$. Le runtime constitué depuis cette
exécution calcule deux terminaux. La même exécution dérive canoniquement son
contrat, le paquet des deux runs et de leur collision opérationnelle, puis un
séparateur fini entre les terminaux ainsi produits. Il s'ensuit que, pour tout
type mémoire et toute réalisation exacte du contrat, leurs images mémoire sont
distinctes. Cette séparation n'est pas un champ de la définition d'exactitude :
elle est dérivée des témoins d'extrémité certifiés, des readouts terminaux et
de leurs lois de préservation. En parallèle, un
regroupement opérationnel à mémoire terminale commune est effectivement
construit lorsqu'on n'exige pas l'exactitude globale des readouts, tandis
qu'une réalisation exacte identité garantit la non-vacuité du résultat.
Enfin, une progression gödelienne finie retient chaque exécution, étend
exactement l'histoire par la phrase produite, puis renouvelle au nouvel
étage le contrat et son obligation universelle. Le développement construit
ainsi, pour toute hauteur naturelle finie, une chaîne dans laquelle chaque
exécution retenue engendre une loi universelle sur toutes ses réalisations
exactes admissibles. Son objet exact est la dérivation, propre à chaque
exécution, d'un contrat et de son obligation, ainsi que leur renouvellement
fini le long des exécutions effectivement produites.
\end{abstract}

\noindent\textbf{Objet central.}
\[
\boxed{\MainAxis}
\]
\noindent\textbf{Formalisation de référence.}
\href{https://github.com/JohnDoe-collab-stack/relational-perimeter}{\code{relational-perimeter}}\\
Révision \code{1e5aeda745559d62a3819bb3213da7fce3c31ad1}.

\noindent\textbf{Claims suivis.}
\begin{itemize}[leftmargin=1.25em,topsep=2pt,itemsep=0pt]
\item \code{GODEL_OPERATIONAL_RUNTIME_CLOSURE_OBSTRUCTION}
\item \code{GODEL_CONSTITUTIVE_THEORY_PROGRESSION}
\end{itemize}

\begin{center}
\fbox{\begin{minipage}{0.91\linewidth}
\textbf{Contrat constitué par une exécution $e$.}
\[
C(e)=\bigl(S_e,I,E,O,\operatorname{next},\operatorname{event},
  \operatorname{read},\operatorname{Allow},\operatorname{decision}\bigr),
\]
avec $O=\mathrm{Type}$ et, pour les deux positions issues du même témoin
gödelien produit,
\[
\operatorname{read}_{C(e)}(t^{\mathrm v})=\Truth(G),\qquad
\operatorname{read}_{C(e)}(t^{\mathrm p})=\Proof_T(G).
\]
Le contrat est fixé avant tout choix de type mémoire et de réalisation.
\end{minipage}}
\end{center}

\section{Changement d'objet de démonstration}

La production d'une phrase gödelienne $G$ certifiée vraie et non prouvable
dans une théorie $T$ demeure un ingrédient indispensable. Cependant, le
résultat présenté ici ne s'arrête pas à ces deux propriétés. L'exécution
qui produit les témoins sert à \emph{constituer un contrat} $C$ ; le contrat
fixe des transitions, des événements, des autorisations et des readouts ;
la distinction attestée par le témoin devient alors une propriété nécessaire
de toute réalisation qui préserve exactement ce contrat.

On passe ainsi d'une affirmation sur le régime de preuve de $T$ à un énoncé
universel portant sur les représentations exactes admissibles de $C$.
L'universalité ne signifie pas que toutes les représentations arbitraires
sont admissibles : elle est quantifiée \emph{sur celles qui réalisent
exactement le contrat de l'exécution}.

\paragraph{Énoncé directeur.}
À chaque étage fini $n$ de la progression considérée, notons $e_n$
l'exécution gödelienne retenue, $C_n$ son contrat,
$t_n^{\mathrm v}$ et $t_n^{\mathrm p}$ ses terminaux de vérité et de preuve,
et $q_A$ la projection mémoire d'une réalisation $A$. Alors :
\begin{equation}\label{eq:main}
\boxed{\begin{gathered}
C_n=\operatorname{godelClosureContract}(e_n),\\[3pt]
\forall M,\;\forall A:\ExactRealization(C_n,M),\\[2pt]
q_A(t_n^{\mathrm v})\neq q_A(t_n^{\mathrm p}).
\end{gathered}}
\end{equation}
La différence des mémoires est une \emph{conclusion} et non une clause
constitutive de $\ExactRealization$.

\section{Constitution du contrat depuis l'exécution}
\label{sec:contract}

\subsection{Données exécutées et témoins}

Soit $T$ un étage doté de la certification requise par l'interface formelle.
La production gödelienne exécutée $e$ contient une phrase $G$ et des
habitants des types
\begin{equation}\label{eq:witnesses}
\Truth(G),\qquad \Proof_T(G)\longrightarrow\Empty.
\end{equation}
Ici $\Proof_T(G)$ désigne le \emph{type de preuve} relatif au régime $T$ ;
le second terme ne constitue pas une preuve de $G$ dans $T$, mais une
fonction qui rejetterait tout habitant d'un tel type.

L'exécution sert également de paramètre à la constitution des deux pôles du
runtime, de ses sources et de ses runs autorisés de fermeture. Ces runs
calculent les terminaux $t^{\mathrm v}$ et $t^{\mathrm p}$, auxquels sont
associés les readouts
\begin{equation}\label{eq:readouts}
\operatorname{read}_C(t^{\mathrm v})=\Truth(G),\qquad
\operatorname{read}_C(t^{\mathrm p})=\Proof_T(G).
\end{equation}
Les égalités concernent des valeurs du type d'observation, qui est ici un
univers de types Lean. Elles ne comparent pas directement des valeurs de
vérité booléennes.

\subsection{Le contrat opérationnel fixé avant toute réalisation}

Dans \SourceAt{RelationalPerimeter/Constitution/Continuation/Behavior.lean}{9-L14}{Behavior.lean}
(lignes 9--14), un \code{FutureContract} comporte, pour des types
$S$ (sources), $I$ (requêtes), $E$ (événements) et $O$ (observations),
les opérations
\[
\begin{aligned}
\operatorname{next}&:S\times I\to S, &
\operatorname{event}&:S\times I\to E,\\
\operatorname{read}&:S\to O, &
\operatorname{Allow}&:S\times I\to\mathrm{Type},
\end{aligned}
\]
ainsi qu'une décision effective d'autorisation ou de rejet pour chaque
couple $(s,i)$. Le contrat prévoit les observations même lorsqu'une requête
est refusée.

Pour l'exécution gödelienne, la définition
\code{godelClosureContract e} de
\SourceAt{RelationalPerimeter/Constitution/Godel/OperationalClosure.lean}{150-L157}{OperationalClosure.lean}
(lignes 150--157) fixe les cinq composantes du contrat : \code{next},
\code{event}, \code{read}, \code{Allow} et \code{decision}. Son
readout est déterminé par le pôle de la source :
\code{Truth(G)} pour le pôle gauche et \code{Proof(G)} pour le pôle droit.
Aucun type mémoire réduit ni aucune réalisation exacte n'entre dans
la définition de ce contrat.

\subsection{Sens exact de la réalisation}

La structure \code{ExactRealization}, définie dans
\SourceAt{RelationalPerimeter/Constitution/Continuation/Minimality.lean}{9-L23}{Minimality.lean}
(lignes 9--23), donne un contrat réduit sur un type mémoire $M$ et une
projection $q_A:S\to M$. Elle impose : préservation bilatérale des
preuves d'autorisation (avec les lois d'aller-retour), commutation avec
les transitions, égalité des événements et exactitude des readouts.
En particulier :
\begin{equation}\label{eq:exactness}
\begin{aligned}
q_A(\operatorname{next}_C(s,i))
  &=\operatorname{next}_{A}(q_A(s),i),\\
\operatorname{event}_C(s,i)
  &=\operatorname{event}_{A}(q_A(s),i),\\
\operatorname{read}_C(s)
  &=\operatorname{read}_{A}(q_A(s)).
\end{aligned}
\end{equation}
Cette structure ne contient aucun champ de la forme
$q_A(t^{\mathrm v})\neq q_A(t^{\mathrm p})$. La préservation est fixée
génériquement ; la nécessité de préserver \emph{cette distinction précise}
provient de l'exécution gödelienne.

\subsection{Paquet constitutif canoniquement dérivé}

Il est utile de concentrer les données produites par une même exécution
sans les transformer en nouvelle primitive. À toute exécution gödelienne
certifiée $e$ sont canoniquement dérivés
\begin{equation}\label{eq:constitutive-package}
e\longmapsto\bigl(C(e),\Gamma_e,\sigma_e\bigr).
\end{equation}
Cette écriture résume une chaîne de constructions déjà séparées dans la
formalisation ; elle ne désigne pas une structure Lean supplémentaire.
Le paquet $\Gamma_e$ retient les deux runs autorisés non vides, leurs
événements et leurs terminaux effectivement calculés :
\[
\Gamma_e=
\bigl(r_e^{\mathrm v},r_e^{\mathrm p},
      t_e^{\mathrm v},t_e^{\mathrm p},\kappa_e\bigr),
\]
où $\kappa_e$ certifie la collision des projections terminales dans la
réalisation opérationnelle sans exactitude globale des readouts. Cette
collision n'identifie ni les sources ni leurs profils constitutifs.

Le dernier composant est construit sur les mêmes terminaux :
\begin{equation}\label{eq:sigma-package}
\sigma_e:
\mathrm{FutureSeparator}
  \bigl(C(e),t_e^{\mathrm v},t_e^{\mathrm p}\bigr),
\qquad \sigma_e.\mathrm{requests}=[].
\end{equation}
Il dépend des témoins et des runs de $e$, mais d'aucun choix de type mémoire
$M$ ni de réalisation exacte $A$. Le contrat fixe donc les règles de
préservation, tandis que $\sigma_e$ détermine la distinction particulière
que ces règles obligent toute réalisation exacte à conserver. On obtient :
\begin{equation}\label{eq:package-obligation}
\boxed{
\sigma_e\Longrightarrow
\forall M,\ \forall A:\ExactRealization(C(e),M),\quad
q_A(t_e^{\mathrm v})\neq q_A(t_e^{\mathrm p}).}
\end{equation}
L'ordre des dépendances est ainsi explicite : l'exécution construit les
runs et leurs terminaux ; les témoins exécutés certifient $\sigma_e$ ;
l'exactitude rend ensuite la séparation obligatoire pour une réalisation
arbitraire du contrat déjà fixé.

\subsection{Deux conséquences distinctes du même contrat}

Le contrat constitué alimente deux chemins formels qu'il ne faut pas
confondre. Le premier impose une obligation universelle à toute réalisation
exacte. Le second part d'un regroupement opérationnel positif et localise
l'incompatibilité d'une préservation dirigée supplémentaire.

\begin{center}
\renewcommand{\arraystretch}{1.25}
\begin{tabular}{>{\centering\arraybackslash}p{0.44\linewidth}
                >{\centering\arraybackslash}p{0.44\linewidth}}
\textbf{Chemin A : séparation universelle}
& \textbf{Chemin B : collision et transport dirigé} \\
paquet $(C(e),\Gamma_e,\sigma_e)$ & paquet $(C(e),\Gamma_e)$ \\
$\downarrow$ & $\downarrow$ \\
séparateur fini $\sigma_e$ avec requête $[]$ & collision opérationnelle $\kappa_e$ \\
$\downarrow$ & $\downarrow$ \\
$\forall A:\ExactRealization(C(e),M)$,
$q_A(t^{\mathrm v})\neq q_A(t^{\mathrm p})$
& collision mémoire et préservation dirigée supplémentaire \\
& $\downarrow$ \\
& transport dirigé $\Truth(G)\to\Proof_T(G)$ \\
& $\downarrow$ avec \code{trueWitness} et \code{unprovable} \\
& $\Empty$
\end{tabular}
\end{center}

Dans le chemin B, la préservation dirigée n'est pas un champ de
$C(e)$ : elle est l'exigence additionnelle dont la compatibilité avec la
collision calculée est testée.

\section{Du témoin gödelien au séparateur fini}
\label{sec:separator}

\begin{lemme}[Incompatibilité des readouts]\label{lem:readout}
Les témoins de \eqref{eq:witnesses} établissent, comme inégalité de types
utilisés par le readout,
\[
\Truth(G)\neq\Proof_T(G).
\]
\end{lemme}
\begin{proof}
Supposons une égalité de ces types. Le transport du témoin
$\Truth(G)$ le long de l'égalité donne un habitant de $\Proof_T(G)$.
L'application du témoin $\Proof_T(G)\to\Empty$ fournit un habitant de
$\Empty$, contradiction. C'est la forme de
\code{godelWitnessReadouts\_ne},
\SourceAt{RelationalPerimeter/Constitution/Godel/OperationalClosure.lean}{636-L641}{OperationalClosure.lean},
lignes 636--641.
\end{proof}

\paragraph{Distinction de deux exécutions de requêtes.}
Les terminaux sont \emph{d'abord calculés} par des runs autorisés de fermeture
portant chacun la requête non vide \code{[close]}. Une fois ces terminaux
obtenus, le séparateur les observe avec la liste de requêtes vide
\code{[]}. Cette seconde liste n'est pas une exécution vide qui aurait
fabriqué les terminaux : elle observe leurs readouts déjà produits.

Le contrat possède une fonction $\operatorname{outcome}_C(s,w)$ qui
assemble readout, autorisations et événements le long de la liste $w$.
Par définition,
\[
\operatorname{outcome}_C(s,[]) =
  \operatorname{stop}(\operatorname{read}_C(s)).
\]
Le lemme précédent et \eqref{eq:readouts} donnent donc :
\begin{equation}\label{eq:separator}
\operatorname{outcome}_{C}(t^{\mathrm v},[])
\neq\operatorname{outcome}_{C}(t^{\mathrm p},[]).
\end{equation}
L'objet \code{godelTerminalSeparator e} construit précisément ce
\code{FutureSeparator}, avec le champ \code{requests := []}
(\SourceAt{RelationalPerimeter/Constitution/Godel/OperationalClosure.lean}{642-L667}{OperationalClosure.lean},
lignes 642--667).

\section{Obligation universelle de réalisation exacte}
\label{sec:theorem}

\begin{lemme}[Principe générique de conservation]\label{lem:generic}
Pour un contrat $C$, une réalisation exacte $A$ de $C$ sur un type mémoire
$M$ et deux sources $s,t$, on a
\[
q_A(s)=q_A(t)\quad\Longrightarrow\quad
\FutureEquivalent(C,s,t).
\]
Ainsi, tout séparateur fini de $s$ et $t$ interdit $q_A(s)=q_A(t)$.
\end{lemme}
\begin{proof}
La préservation de la transition, de l'autorisation, de l'événement et du
readout donne, par induction sur les listes de requêtes, l'égalité des
outcomes du contrat et du contrat réduit. Deux mémoires projetées égales
ont alors les mêmes outcomes réduits pour toute liste ; par transport,
les outcomes sources sont égaux. Un séparateur fournit une liste pour
laquelle cette égalité est impossible. Ces étapes sont les théorèmes
\code{ExactRealization.outcome\_exact},
\code{equal\_memory\_same\_futures}, puis
\code{distinct\_memory\_of\_separator} dans
\SourceAt{RelationalPerimeter/Constitution/Continuation/Minimality.lean}{56-L75}{Minimality.lean}
et \SourceAt{RelationalPerimeter/Constitution/Continuation/ExactSeparation.lean}{20-L27}{ExactSeparation.lean}.
\end{proof}

\begin{theoreme}[Séparation universelle engendrée par l'exécution gödelienne]
\label{thm:universal}
Soit $e$ une exécution gödelienne certifiée et $C(e)$ le contrat
\code{godelClosureContract e} qu'elle constitue. Soient
$t^{\mathrm v}(e)$ et $t^{\mathrm p}(e)$ les terminaux effectivement
calculés par le regroupement opérationnel de cette exécution. Alors
\begin{equation}\label{eq:universal}
\boxed{\begin{gathered}
\forall M,\quad \forall A:\ExactRealization(C(e),M),\\
q_A(t^{\mathrm v}(e))\neq q_A(t^{\mathrm p}(e)).
\end{gathered}}
\end{equation}
\end{theoreme}
\begin{proof}
L'exécution produit les témoins de \eqref{eq:witnesses} et ses deux runs
calculent les terminaux. Le lemme~\ref{lem:readout} fournit la différence
de readouts ; la requête vide fournit le séparateur fini
\eqref{eq:separator}. Pour une réalisation exacte $A$ arbitraire, le
lemme~\ref{lem:generic} interdit l'égalité de leurs projections mémoire.
Aucun choix de $M$ ou de $A$ n'intervient dans la construction du contrat
ni dans celle du séparateur. L'énoncé se conclut donc pour toutes les
réalisations exactes admissibles. En Lean, le théorème
\begin{center}
\code{godelEveryExactRealizationSeparatesTerminals}
\end{center}
est obtenu par
application directe de
\code{realization.distinct\_memory\_of\_separator}
à \code{godelTerminalSeparator execution}
(\SourceAt{RelationalPerimeter/Constitution/Godel/OperationalClosure.lean}{669-L676}{OperationalClosure.lean},
lignes 669--676).
\end{proof}

\begin{remarque}[Nature de l'universalité]
Le contrat $C(e)$ \emph{dépend} de l'exécution $e$. Ce qui est indépendant
d'une réalisation particulière n'est ni l'existence de cette exécution
ni le choix de son contrat, mais la preuve de la contrainte
\eqref{eq:universal}. La quantification porte sur les réalisations exactes
de $C(e)$, et non sur tous les modèles possibles de Gödel.
\end{remarque}

\begin{remarque}[Force minimale de l'obstruction]
Pour ces terminaux, la liste séparatrice est vide : le désaccord est
immédiatement visible dans les readouts. L'argument particulier peut donc
être lu directement via \code{read\_exact}. Le lemme générique de
conservation des futurs est plus large que ce qui est nécessaire pour
ce seul séparateur. Cette observation ne modifie ni la constitution du
contrat ni la portée universelle du théorème.
\end{remarque}

\section{Frontière exacte : regroupement, conservation, existence}
\label{sec:frontier}

Le théorème précédent ne dit pas qu'aucun regroupement des deux branches
n'est exécutable. Le code sépare précisément trois affirmations.

\begin{proposition}[Regroupement opérationnel sans exactitude globale]
Une projection opérationnelle conserve les transitions, les événements
et les autorisations tout en calculant une mémoire terminale commune après
les deux runs de fermeture. Elle n'assume pas l'exactitude globale des
readouts.
\end{proposition}
\begin{proof}[Construction]
La mémoire opérationnelle comporte des états actifs retenant le pôle et
des états \code{grouped} qui ne le retiennent plus. La transition de
fermeture produit un état \code{grouped} commun. Les lois opérationnelles
sont établies dans \code{godelOperationalRealizationWithoutRead} ;
\code{godelOperationalGroupingWithoutRead} exécute les deux runs et
démontre leur collision calculée (\SourceAt{RelationalPerimeter/Constitution/Godel/OperationalClosure.lean}{589-L602}{OperationalClosure.lean},
lignes 165--272 et 358--372, 589--602). Le contrat réduit lit ici une
observation constante \code{PUnit} : il ne prétend pas reproduire les deux
readouts sources.
\end{proof}

\begin{proposition}[Rejet de la conservation dirigée en présence de collision]
Ajouter à ce regroupement la conservation dirigée permettant de passer du
readout de vérité au readout de preuve produirait
$\Truth(G)\to\Proof_T(G)$, puis $\Empty$.
\end{proposition}
\begin{proof}
Le regroupement calcule l'égalité des mémoires terminales. Une
\begin{center}
\code{DirectedTerminalReadoutPreservation}
\end{center}
transporterait alors un
habitant du readout de vérité vers le readout de preuve. Appliquée au
témoin vrai et suivie du certificat de non-prouvabilité, cette conversion
produit $\Empty$. C'est la chaîne
\begin{center}
\code{godelComputedGroupingToProducedCompleteness}\\[2pt]
$\downarrow$\\[2pt]
\code{rejectGodelComputedGroupingToEmpty}
\end{center}
Elle est distincte de la
preuve du théorème~\ref{thm:universal}.
\end{proof}

\begin{proposition}[Non-vacuité de l'exactitude]
Pour tout contrat $C$ de type \code{FutureContract}, il existe une
réalisation exacte identité, dont le type mémoire est celui des sources.
En particulier, le contrat gödelien possède une réalisation exacte.
\end{proposition}
\begin{proof}[Construction]
La définition \code{identityRealization} pose
$M=S$, $q_A=\mathrm{id}$ et reprend le contrat comme contrat réduit.
Les lois d'exactitude se réduisent à des égalités réflexives ; voir
\SourceAt{RelationalPerimeter/Constitution/Continuation/CommonMemoryGrouping.lean}{10-L19}{CommonMemoryGrouping.lean},
lignes 10--19. La spécialisation est \code{godelIdentityRealization},
puis \code{heytingArithmeticIdentityRuntime} pour l'instance HA.
\end{proof}

Ces trois résultats fixent la frontière : \emph{fusion opérationnelle
constructible sans exactitude globale des readouts ; fusion des terminaux
interdite dans toute réalisation exacte du contrat ; existence effective
d'au moins une telle réalisation exacte}. La non-vacuité est globale : elle
n'implique pas que chaque type mémoire arbitraire admette une réalisation
exacte.

\section{Renouvellement constitutif à chaque étage fini}
\label{sec:progression}

Un \code{SoundGodelTheoryStage} conserve une histoire $\mathcal H_n$, le
régime arithmétique $T_n$ qu'elle présente, sa certification,
une exécution gödelienne maîtresse et une preuve d'identité de cette exécution
avec la production canonique du régime. La procédure
\code{executeGodelTheoryProgressionFrom} consomme cette exécution retenue.
Elle ajoute \emph{l'occurrence exacte} de la phrase qu'elle a produite :
\begin{equation}\label{eq:successor}
\mathcal H_{n+1}
  =\operatorname{extend}\bigl(\mathcal H_n,
    \operatorname{occurrence}(G_n)\bigr).
\end{equation}
Cette écriture désigne une extension constitutive de l'histoire par
l'occurrence effectivement produite. Elle ne désigne ni une union
ensembliste primitive ni l'ajout d'une phrase reconstruite indépendamment.
La certification s'appuie sur le témoin de vérité de la production source.
L'occurrence nouvellement ajoutée fournit une preuve vérifiée de $G_n$
dans $T_{n+1}$. Cette preuve est construite \emph{par l'occurrence ajoutée}
: la fonction Lean nommée \code{oldTruthToTargetProof} ignore son argument
de vérité et renvoie la preuve cible déjà construite. Il ne s'agit donc
pas d'un procédé général transformant la vérité en preuve.

La même étape exécute une \emph{nouvelle} production $G_{n+1}$ relativement
au régime étendu. On dispose alors des données suivantes :
\[
\begin{gathered}
\Proof_{T_n}(G_n)\to\Empty,
\qquad \Proof_{T_{n+1}}(G_n),\\
\Truth(G_{n+1}),\qquad
\Proof_{T_{n+1}}(G_{n+1})\to\Empty,\\
G_n\neq G_{n+1}.
\end{gathered}
\]
L'inégalité des phrases est démontrée : les identifier transporterait
la preuve construite de $G_n$ au nouvel énoncé non prouvable
$G_{n+1}$ (\SourceAt{RelationalPerimeter/Instances/HeytingArithmetic/GodelTheoryProgression.lean}{175-L207}{GodelTheoryProgression.lean},
lignes 175--207).

\begin{theoreme}[Renouvellement fini de l'obligation]\label{thm:finite}
Pour tout étage certifié $S_0$ et toute hauteur naturelle $n$, l'étage
$S_n=\operatorname{advance}(S_0,n)$ est constitué par itération finie.
En désignant par $e_n$ l'exécution qu'il retient, par $C_n=C(e_n)$ son
contrat et par $t_n^{\mathrm v}$, $t_n^{\mathrm p}$ ses terminaux calculés,
la conclusion \eqref{eq:main} vaut à cet étage. Cette conclusion vaut en
particulier pour la progression issue de l'étage initial HA.
\end{theoreme}
\begin{proof}
La construction \code{advance} est récursive sur $n$ ; la construction
\code{executedTrace} retient en outre chaque étape, son exécution cible
et son histoire. Son cas successeur consomme le terminal effectivement
produit par l'étape précédente, plutôt que de reconstruire une histoire
équivalente à partir d'un simple indice. La propriété universelle
\code{everyExactRealizationSeparatesCurrentTerminals} s'applique à
chaque \code{SoundGodelTheoryStage}, donc à l'étage obtenu pour la hauteur
$n$ (\SourceAt{RelationalPerimeter/Instances/HeytingArithmetic/GodelTheoryProgression.lean}{291-L448}{GodelTheoryProgression.lean},
lignes 291--448).
\end{proof}

L'extension de la théorie et la réalisation du contrat restent deux
opérations distinctes. Aucune réalisation exacte commune n'est construite
entre $C_n$ et $C_{n+1}$ ; les frontières de preuve et les obligations sont
réévaluées \emph{relativement au contrat propre de chaque étage}.

La trace exécutée conserve les étapes et les exécutions effectivement
produites. Les obligations universelles ne sont pas stockées comme des champs
de l'histoire ou de la trace. Elles sont dérivables à chaque étage retenu en
appliquant \code{everyExactRealizationSeparatesCurrentTerminals} à l'exécution
propre de cet étage. À chaque étage de la trace finie, une obligation
renouvelée, relative à son propre contrat, est donc dérivable.

\section{Contribution et portée exacte du résultat}

La contribution formellement présentée est la chaîne constitutive complète
\[
\boxed{\begin{gathered}
e\longmapsto\bigl(C(e),\Gamma_e,\sigma_e\bigr)\\[3pt]
\downarrow\\[3pt]
\forall M,\;\forall A:\ExactRealization(C(e),M),\\[2pt]
q_A(t_e^{\mathrm v})\neq q_A(t_e^{\mathrm p}).
\end{gathered}}
\]
Elle est réalisée par la production gödelienne, son séparateur fini, le
théorème universel de séparation et leur renouvellement à toute hauteur finie.

\paragraph{Découverte présentée.}
Une exécution gödelienne certifiée ne fournit pas seulement une phrase et
des propriétés relatives à sa prouvabilité. À partir du même témoin exécuté,
elle constitue son contrat de futurs, produit deux runs et leurs terminaux,
construit une collision opérationnelle sous le régime qui n'impose pas
l'exactitude globale des readouts, puis construit sur ces mêmes terminaux un
séparateur fini. Cette donnée particulière engendre alors une conséquence
indépendante de tout choix ultérieur de mémoire : chaque réalisation exacte
admissible du contrat doit maintenir distinctes les projections des deux
terminaux. La progression conserve les exécutions produites et renouvelle
cette génération d'une obligation universelle à chaque étage naturel fini.

Le passage décisif n'est pas la seule inégalité finale entre deux mémoires.
À partir d'une exécution singulière sont construits le contrat, les runs, les
terminaux et le séparateur auxquels s'applique ensuite la loi générique
d'exactitude. La distinction particulière imposée aux réalisations n'est donc
ni un champ de \code{ExactRealization}, ni une séparation choisie avant
l'exécution. Elle est engendrée par les objets effectivement produits, puis
devient universelle relativement au contrat ainsi constitué. Le résultat
transforme donc les témoins gödeliens certifiés en une obligation
représentationnelle portant sur toute réalisation exacte admissible.

La contribution ne réside donc ni dans le seul lemme abstrait affirmant
qu'une observation préservée doit rester distinguable, ni dans la seule
production d'une phrase gödelienne. Elle réside dans leur articulation
constitutive complète : les données auxquelles s'applique l'obligation
universelle sont précisément celles qu'a produites et retenues l'exécution
gödelienne qui constitue le contrat.

Le résultat peut être résumé en une phrase :
\begin{quote}
Une exécution gödelienne certifiée dérive canoniquement un contrat, deux runs
et leurs terminaux, une collision opérationnelle ainsi qu'un séparateur fini
de ces mêmes terminaux. Sans postuler d'inégalité entre leurs mémoires
projetées, ce séparateur impose leur séparation dans toute réalisation exacte
admissible ; la progression des théories renouvelle cette obligation à chaque
étage fini constitué.
\end{quote}

\paragraph{Portée exacte.}
Le théorème porte sur les contrats constitués par les exécutions certifiées
de la formalisation et sur toutes leurs réalisations exactes admissibles.
La progression est construite pour chaque hauteur naturelle finie ; elle
retient les occurrences et les exécutions effectivement produites, puis
réévalue l'obligation relativement au contrat propre du nouvel étage. Son
objet ne comprend pas d'étage limite, d'union transfinie, de théorie finale
complète ni de réalisation transversale entre les contrats successifs. La
séparation démontrée concerne la non-prouvabilité certifiée de la phrase
produite ; elle ne remplace pas une preuve distincte de non-réfutabilité.

\paragraph{Positionnement scientifique.}
La diagonalisation exécutée, la séparation observationnelle, les contrats
indépendants de leurs implémentations et les progressions de théories ont
chacun des antécédents établis. La contribution revendiquée ici n'est aucun
de ces ingrédients isolés. Dans le corpus primaire comparatif examiné, aucun
antécédent n'a été identifié qui construise, à partir du même témoin
gödelien exécuté, la chaîne entière
\[
e\longmapsto\bigl(C(e),\Gamma_e,\sigma_e\bigr)
\longmapsto
\forall M\,\forall A:\ExactRealization(C(e),M),\quad
q_A(t_e^{\mathrm v})\neq q_A(t_e^{\mathrm p}),
\]
qui établisse en même temps la possibilité du regroupement opérationnel sans
exactitude globale, la non-vacuité de la quantification par une réalisation
exacte positive, puis renouvelle cette obligation sur la trace finie
effectivement produite. Ce positionnement localise la contribution ; il ne
remplace ni la chaîne formelle ni l'axe d'exposition qui part de l'exécution.

\section{Noyau Lean et traçabilité des dépendances}
\label{sec:lean}

La formule centrale rassemble quatre constructions formelles distinctes :
\begin{center}
\renewcommand{\arraystretch}{1.2}
\begin{tabular}{>{\raggedright\arraybackslash}p{0.28\linewidth}
                >{\raggedright\arraybackslash}p{0.61\linewidth}}
\toprule
\textbf{Notation du manuscrit} & \textbf{Construction Lean} \\
\midrule
$C(e)$ & \code{godelClosureContract execution} \\
$\Gamma_e$ & \code{godelOperationalGroupingWithoutRead execution} \\
$\sigma_e$ & \code{godelTerminalSeparator execution} \\
obligation universelle & \code{godelEveryExactRealizationSeparatesTerminals} \\
\bottomrule
\end{tabular}
\end{center}

Le triplet $(C(e),\Gamma_e,\sigma_e)$ est une notation documentaire pour
cette composition de constructions existantes. Il ne désigne pas une
structure Lean supplémentaire.

Le point de passage qui rend visible le changement d'objet est le théorème
\begin{center}
\code{godelEveryExactRealizationSeparatesTerminals}
\end{center}
de
\SourceAt{RelationalPerimeter/Constitution/Godel/OperationalClosure.lean}{669-L676}{OperationalClosure.lean}
(lignes 669--676). Son terme de preuve est exactement :
\begin{Verbatim}[fontsize=\small]
realization.distinct_memory_of_separator (godelTerminalSeparator execution)
\end{Verbatim}
La progression aboutit au théorème suivant, valable pour un étage certifié
quelconque :
\begin{center}
\code{everyExactRealizationSeparatesCurrentTerminals}
\end{center}
Son hypothèse \code{realization} est quantifiée sur toute réalisation exacte
du contrat de l'exécution retenue par l'étage.

\paragraph{Dépendances décisives, dans leur ordre logique.}
\begin{enumerate}[leftmargin=*,itemsep=3pt]
\item \textbf{Constitution.} \code{godelClosureContract} est défini à partir de l'exécution, sans mémoire choisie.
\item \textbf{Production de $\Gamma_e$.} \code{godelOperationalGroupingWithoutRead} exécute les deux runs non vides, retient leurs événements et leurs terminaux, puis construit leur collision opérationnelle sans exactitude globale des readouts.
\item \textbf{Construction de $\sigma_e$.} \code{godelTerminalSeparator} réutilise les terminaux de ces runs, fixe la requête \code{[]} et prouve le désaccord de leurs outcomes à partir des témoins de la même exécution.
\item \textbf{Passage au général.} \code{distinct_memory_of_separator} transforme le séparateur en contrainte sur chaque réalisation exacte, sans choisir la mémoire ni la réalisation lors de la construction de $\sigma_e$.
\item \textbf{Frontière.} La collision de $\Gamma_e$ est positive dans la réalisation opérationnelle sans readout exact, tandis que \code{identityRealization} construit une réalisation exacte qui conserve les terminaux distincts.
\item \textbf{Progression.} \code{executedTrace} conserve les étapes ; l'obligation est portée à un étage quelconque par le théorème indiqué ci-dessus.
\end{enumerate}

\paragraph{Statut de la vérification.}
Les définitions et les corps de preuve cités ont été confrontés aux fichiers
publics du dépôt à la révision indiquée en première page. Les sources
comportent des directives \texttt{\#print axioms} pour les déclarations
centrales. Leur présence ne vaut pas, à elle seule, rapport d'audit des
résultats ; ce manuscrit ne prétend pas avoir exécuté indépendamment
\texttt{lake build} ni reproduit les sorties de l'audit axiomatique.

\paragraph{Statut documentaire.}
Ce texte réorganise, sans augmenter la portée des résultats, le théorème-cadre
\Source{docs/theoreme-cadre-progression-godelienne-separation-exacte.fr.tex}{LaTeX}
et la présentation canonique
\href{https://github.com/JohnDoe-collab-stack/relational-perimeter/blob/1e5aeda745559d62a3819bb3213da7fce3c31ad1/docs/progression-constitutive-theories-godeliennes.fr.md}{\emph{Progression constitutive des théories et obstructions gödeliennes}}.
Il explicite avant tout l'axe exécution--contrat--obligation universelle,
qui motive le présent ordre d'exposition.

\section{Conclusion}

Le déplacement établi par la formalisation part de l'exécution elle-même.
Une production gödelienne certifiée constitue les objets, le contrat et le
témoin fini qui rendent ensuite démontrable une obligation portant sur toute
réalisation exacte admissible, quel que soit son type mémoire. L'universel
n'est donc ni ajouté extérieurement à l'exécution ni supposé dans la
définition de l'exactitude : il est engendré par la conservation exacte des
distinctions que cette exécution particulière a effectivement produites.

La progression finie répète ce passage sans reconstruire indépendamment ses
entrées : chaque nouvel étage consomme l'histoire et l'exécution retenues,
constitue son propre contrat, puis renouvelle l'obligation universelle. Le
résultat central demeure ainsi celui annoncé dès la première page :
\[
\boxed{\MainAxis}
\]

\end{document}
